%% kapa: exam=1 round=19 year=2008 period=1 form=A subject=civil_law subject_ko=민법 seq=32 number=32
%% source: Q-Net 감정평가사 기출문제 (cst00309) artlSeq=1001962 file=extracted/1008/1001962/2008년 제19회 감정평가사 1차시험 기출문제/1교시.hwp
%% conversion: kordoc 4.24.51 -> tools/gen1.py (auto); review=proofread-sonnet; flags=none
\begin{question}{32}
건물의 구분소유에 관한 설명으로 옳은 것은?
\begin{choices}
  \item 공용부분은 공유자 전원에게 공동으로 귀속되므로 일부 구분소유자에게만 귀속될 수는 없다.
  \item 구분소유자간의 규약으로 건물의 공용부분이 성립할 수 없다.
  \item 전유부분은 구조상의 독립성과 기능상의 독립성을 모두 가져야 한다.
  \item 공유자의 공용부분에 대한 지분은 균등함이 원칙이다.
  \item 공유자는 그가 가지는 전유부분과 분리하여 공용부분에 대한 지분을 처분할 수 있다.
\end{choices}
\end{question}
\answer{3}
